\section{Discussion}
\label{sec:discussion}

\paragraph{What the results mean.}
Mixed-resolution input is not a cheaper version of sharp input for these models; it is input they handle poorly. For \pithree and \dathree the cause lies mainly in the mismatch of sizes rather than in missing detail, and it is not only a matter of cross-view positions, since \dathree has none. Because the problem sits at the input, the most effective training-free fix is also at the input: give every view the same size. A natural next step is training with mixed-resolution views. Our results suggest what such training would need to cover: views of different token counts within one input, not only lower resolution overall.

\paragraph{Failed predictions.}
Each study's predictions were fixed before it ran, and we report those that failed. \dathree was predicted to barely degrade and to gain nothing from alignment (\cref{sec:diagnosis}). A calibrated attention correction was predicted to improve on the uncalibrated one, and a correction strength chosen on held-out data was predicted to transfer to the other datasets (\cref{sec:fixes}). A corrected mix was predicted to match or beat the best uniform input; it was below it in all 8 cells, 3 of them significantly, which is why this paper studies why mixing fails rather than how to make it win. The mix was predicted to gain at low budgets (\cref{sec:budget}). VGGT's gap between true and upsampled input was predicted to shrink at the higher token level, and equal detail was predicted to help it there (\cref{sec:vggt,sec:fixes}). Shrinking every view was predicted to beat the mix at the higher level for \dathree and VGGT as well as \pithree. Two predictions failed narrowly: twelve views at one size lost to the four sharp views in only 5 of 10 cells for \dathree with graded mixes, and on the replication seed the one-size fix lost one cell for \dathree.

\paragraph{Limitations.}
Our mixes are simulated: small views are downsampled from the same images as the sharp ones. This is what makes a true reference possible, but real mixed captures also differ in sensors, optics and exposure, and ratios beyond 4 or continuous ranges of sizes are untested. We use the published models without retraining. The core diagnosis rests on two frame samples that agree in 188 of 189 cells significant on both; the other studies use the reporting seed only, with the selection seed agreeing in direction. We do not establish why VGGT differs, and its cells at the lower token level sit near its floor. We do not correct for multiple comparisons; no claim rests on a single cell, and our conclusions rest on consistent counts such as 20 of 20. The common-frame scoring was added after the first results were read; it agrees with all-pair scoring. Reconstruction covers fewer settings than pose: ETH3D uses 9 of 11 scenes because two exceed the memory of the fusion step, the studies of splits, ratios, graded mixes, view counts, shrinking and replication are pose only, and apart from the held-out DTU comparison, some reconstruction results come from readouts in which not every job had yet been fused. Cost is measured in tokens, which our timing supports (\cref{sec:setup}).

\paragraph{Conclusion.}
Feed-forward 3D models collapse when the views of one input differ in resolution. For \pithree and \dathree, upsampling the small views restores most of the lost accuracy without adding information, which places the problem in how the views are represented. Corrections inside the model recover little; giving every view one size, at the mix's own cost, recovers most. And when a budget can be spent freely, the benefit often attributed to mixing resolutions turns out to be the benefit of using more views.
